\chapter{Bilan, limites et passation}
\label{ch:bilan}

\questionchapitre{Que reste-t-il, une fois le contrat achevé, et dans quel état ?}

\section{Récapitulatif}

\begin{center}
\begin{tabularx}{\textwidth}{@{}p{3.3cm}X@{}}
\toprule
\textbf{Domaine} & \textbf{Réalisation} \\
\midrule
Données & Chaîne d'ingestion automatisée de quatre jeux de données issus de trois fournisseurs
publics ; entrepôt en trois couches sur base relationnelle étendue au temporel et au
géospatial ; rattachement spatial des stations à leur maille climatique ; enrichissement par le
référentiel hydrogéologique ; environ \num{658} millions de lignes en couche brute \\
Fiabilité & Idempotence des traitements, traçabilité des lignes écartées, contraintes de
concurrence par clés, reprise après interruption, procédures de sauvegarde et de restauration
vérifiées à l'échelle réelle \\
Indicateurs & Indices standardisés de niveau, de débit, de précipitation, de température et de
bilan hydrique ; production nationale sur quatre fenêtres ; protocole de validation à critère
préétabli ; deux décisions méthodologiques établies par la mesure \\
Exploration & Observatoire des nappes et des rivières ; module climatique national avec
historique depuis 1950, épisodes de sécheresse et export \\
Modélisation métier & Laboratoire de modèles à fonction de transfert ; calibration automatique
guidée par la nature de l'aquifère ; évaluation selon quatre critères normalisés ; diagnostics
préalables et postérieurs ; scénarios à bornes physiquement dérivées \\
Apprentissage & Douze architectures de prévision profonde ; optimisation d'hyperparamètres ;
traçabilité des expériences ; suivi d'entraînement en temps réel ; configurations prêtes à
l'emploi par cas d'usage \\
Explicabilité & Attribution multi-méthodes ; contrefactuels sous contrainte physique avec
couplage thermodynamique et validation par modèle indépendant ; détection non supervisée
d'influence anthropique \\
Ingénierie & Deux dépôts conteneurisés, reproductibles et documentés pour la reprise ; environ
\num{1400} contributions ; environnements de développement et de production séparés \\
\bottomrule
\end{tabularx}
\end{center}

\section{Bilan au regard de la mission}

La mission consistait à amorcer, dans le cadre de \textsc{Junon}, les outils et les données sans
lesquels les travaux d'explicabilité d'\textsc{Aida} ne pourraient être menés.

\textbf{Ce qui est atteint} : l'infrastructure complète, de l'ingestion des sources ouvertes
jusqu'aux méthodes d'explication, avec une contribution originale sur les contrefactuels sous
contrainte physique qui répond précisément à la difficulté identifiée avec les hydrogéologues,
à savoir qu'une explication physiquement incohérente n'est pas une explication.

\textbf{Ce qui n'est pas atteint} : l'évaluation quantitative comparée des méthodes, faute de
temps une fois l'infrastructure construite. C'est le regret principal de la période. C'est aussi
le travail le plus immédiatement accessible à qui reprendra ces développements, parce que tout
ce qu'il suppose existe désormais.

\begin{remarque}[title={Un arbitrage à assumer}]
Consacrer l'essentiel d'un postdoctorat à construire l'infrastructure d'un sujet plutôt qu'à
publier sur ce sujet est défavorable à court terme et favorable à long terme. Le jeu de données,
les indicateurs validés et les outils demeurent, et ils rendent possibles des travaux qui ne
l'étaient pas. La question posée au chapitre~\ref{ch:verrou} a reçu une réponse construite
plutôt qu'argumentée.
\end{remarque}

\section{Limites assumées}

Elles sont énumérées sans atténuation, parce qu'un successeur doit les connaître et qu'un
rapport qui n'en comporterait aucune ne serait pas crédible.

\begin{enumerate}
  \item \textbf{Aucun résultat de performance comparée n'est établi.} Les campagnes
  d'entraînement ont validé le dispositif, non classé les méthodes. Les jeux préparés et les
  points de sauvegarde ont été purgés lors de la passation.

  \item \textbf{Une incohérence connue subsiste sur l'évapotranspiration} entre la chaîne de
  grille et la chaîne de station. La seconde conserve la variable d'origine parce qu'elle
  constitue l'entrée de forçage des modèles métier déjà calibrés et que sa modification les
  invaliderait. L'arbitrage est explicite et documenté ; sa résolution suppose une recalibration
  d'ensemble.

  \item \textbf{Les tests ne s'exécutent pas automatiquement} sur la plateforme, la chaîne
  d'intégration de la forge hébergeante ne comportant qu'une étape de construction. Quatre des
  \num{552} tests échouent, pour des raisons documentées et sans rapport avec un défaut
  applicatif, mais la purge des journaux se trouve de fait non testée.

  \item \textbf{Certaines températures agrégées au niveau des stations} sont des extrêmes de
  moyennes journalières et non de véritables extrêmes journaliers ; les exposer supposerait de
  modifier une table partitionnée.

  \item \textbf{La source climatique retenue est un compromis.} ERA5-Land a été préféré à une
  réanalyse plus fine mais non redistribuable, au bénéfice de la reproductibilité et du partage.

  \item \textbf{L'usage reste circonscrit au cercle du projet}, faute d'une politique d'accès
  élargie.

  \item \textbf{Le périmètre couvre la quantité, non la qualité} des eaux, également visée par
  le programme.
\end{enumerate}

\section{État de passation}

Le contrat s'est achevé le 30 juin 2026. Les travaux se sont poursuivis jusqu'au 24 août 2026
dans le seul but de rendre les réalisations exploitables sans leur auteur :

\begin{itemize}
  \item refonte complète de la documentation des deux dépôts, organisée autour d'une
  cartographie des documents et vérifiée par exécution réelle des procédures décrites ;
  \item consignation des pièges susceptibles de produire des résultats faux sans émettre
  d'erreur, développés au chapitre~\ref{ch:entrepot} ;
  \item retrait des éléments sensibles, identifiants, clés d'accès, topologie interne et
  adresses personnelles, et restriction par défaut de l'écoute des services à l'interface
  locale ;
  \item suppression du code mort et des scripts devenus inopérants ;
  \item vérification à l'échelle réelle des procédures de sauvegarde et de restauration ;
  \item explicitation du couplage entre les deux dépôts et de ses conditions de démarrage.
\end{itemize}

\section{Conclusion}

Ces travaux répondent à une question préalable à celle qui les motivait : comment rendre
possible la recherche en explicabilité sur les eaux souterraines françaises.

La réponse tient en une chaîne complète, de quatre jeux de données publics et dispersés
jusqu'à des scénarios contrefactuels énonçables dans le langage des hydrogéologues, dont chaque
maillon est documenté, vérifié et réutilisable indépendamment des autres. Le coût d'entrée à
l'expérimentation, qui se comptait en semaines et se payait autant de fois qu'il y avait
d'équipes, se réduit à une requête.

Ce qui a été construit dépasse le périmètre des eaux souterraines. L'enchaînement des sources
consolidées vers les indicateurs validés, puis vers les modèles et l'explication, est
transposable aux autres compartiments environnementaux visés par le programme. C'est en ce sens
que ces développements constituent moins un aboutissement qu'un point de départ, ce qui était
précisément leur objet.
